% !TeX root = se100_the_ontological_neutrality_theorem.tex

\section{Formal Development of the Ontological Neutrality Theorem}
\label{sec:formal-dev}

Section~\ref{sec:requirements} defined the neutrality
requirements governing a substrate ontology:
interpretive non-commitment and extension stability.
It also introduced the scope assumptions under which those requirements apply:
causal and normative propositions are contestable across admissible frameworks,
while the substrate's framework-invariant referential structure,
including its referential regime,
is assumed to be invariant
across the admissible frameworks under consideration.

This section establishes the relationship between those requirements
and the structure of the substrate layer.
Specifically, it proves that the two neutrality requirements
place a structural constraint on the substrate layer:
under the contestability assumption and
the referential-invariance scope condition,
a substrate satisfies interpretive non-commitment and extension stability
if and only if its foundational layer is
restricted to framework-invariant referential structure and
makes no causal or normative commitments.

The argument is structural and logical rather than empirical.
It follows from the interaction between interpretive disagreement
and substrate-layer assertion.
All necessity and sufficiency claims in this section
are relative to the definitions, assumptions, and scope conditions introduced in
Section~\ref{sec:requirements}.

%-------------------------------------------------------
\subsection{Neutrality as Compatibility Across Admissible Frameworks}
\label{subsec:compatibility}
%-------------------------------------------------------

By definition, a neutral substrate must remain compatible
with multiple admissible interpretive frameworks,
even when those frameworks disagree.
Compatibility here is understood in the sense of extension stability:
the substrate must remain logically consistent
when extended by each admissible interpretive framework,
without revision or retraction of substrate-layer assertions.

Admissible frameworks may disagree about
causal or normative conclusions concerning the same referents
while remaining internally consistent.
For example, an admissible framework may correspond to
a recognized legal jurisdiction, an established causal modeling tradition,
or a coherent institutional or normative charter.
Admissibility is not mutual compatibility.
Two admissible frameworks may contradict one another,
as when two jurisdictions reach opposite legal conclusions
or two causal models attribute responsibility differently.
The set $\mathbb{F}$ is therefore characterized by domain-governed
admissibility and internal consistency,
not by global agreement among its members.

This neutrality requirement immediately constrains
what the substrate may assert as foundational fact.
If the substrate asserts a proposition whose truth depends
on a particular causal or normative interpretation,
then that assertion may conflict with an admissible framework
that rejects the proposition.
Because neutrality requires the substrate to remain consistent
when extended by every admissible framework,
such an assertion violates the neutrality requirements.

This restriction does not prohibit the representation of
causal or normative concepts.
A neutral substrate may record
causal or normative claims as reified assertions
carrying attribution and provenance
without asserting their content as substrate-layer facts.
For example, the substrate may assert
$\mathsf{Asserts}(\mathcal{F}, \varphi)$,
where $\varphi$ is a causal or normative claim,
without asserting $\varphi$ itself.
What neutrality excludes is the substrate committing to
causal or normative propositions about substrate individuals.

%-------------------------------------------------------
\subsection{Framework-Contestability and Substrate Incompatibility}
\label{subsec:general-lemma}
%-------------------------------------------------------

The arguments in the following subsections share a common structure,
stated here as a general lemma.

\begin{lemma}[Framework-Contestability Lemma]
  \label{lem:contestability}
  Let $\mathcal{S}$ be a substrate ontology intended to satisfy the neutrality requirements,
  and let $p$ be a framework-variant proposition
  with respect to $\mathbb{F}$.
  If $\mathcal{S}$ commits to $p$,
  then $\mathcal{S}$ violates both interpretive non-commitment
  and extension stability.
\end{lemma}

\begin{proof}
  Since $p$ is framework-variant with respect to $\mathbb{F}$,
  there exist admissible frameworks
  $\mathcal{F}_1,\mathcal{F}_2 \in \mathbb{F}$
  such that
  $\mathcal{S}\cup\mathcal{F}_1 \vdash p$
  and
  $\mathcal{S}\cup\mathcal{F}_2 \vdash \neg p$.

  Since $\mathcal{S}$ commits to $p$,
  Definition~\ref{def:substrate-layer-commitment}
  gives $\mathcal{S}\vdash p$.
  Thus $\mathcal{S}$ asserts a framework-variant proposition,
  violating interpretive non-commitment
  (Definition~\ref{def:interpretive-non-commitment}).

  Moreover, because $\mathcal{S}\vdash p$,
  every extension of $\mathcal{S}$ preserves $p$.
  In particular,
  $\mathcal{S}\cup\mathcal{F}_2 \vdash p$.
  But by framework-variance,
  $\mathcal{S}\cup\mathcal{F}_2 \vdash \neg p$.
  Hence
  \[
    \mathcal{S}\cup\mathcal{F}_2 \vdash \bot .
  \]
  Thus $\mathcal{S}$ also violates extension stability
  (Definition~\ref{def:extension-stability}).
\end{proof}

The lemma identifies the general failure mode:
a substrate loses neutrality when it commits to
a proposition that admissible frameworks may contest.
The next two subsections instantiate this structure
for normative and causal commitments.

%-------------------------------------------------------
\subsection{Normative Commitments}
\label{subsec:normative}
%-------------------------------------------------------

Normative conclusions are framework-dependent
in the sense of Lemma~\ref{lem:contestability}.
Their truth conditions depend on legal regimes,
institutional authorities, temporal scope, governing rules, and interpretive stance.
It is therefore a routine and expected feature of accountability systems
that admissible frameworks disagree about normative conclusions.

By Assumption~\ref{asm:contestability},
if $\mathcal{S}$ commits to a normative proposition $n$,
then $n$ is not guaranteed invariant across admissible frameworks:
there exists some $\mathcal{F} \in \mathbb{F}$ such that
$\mathcal{S} \cup \mathcal{F} \vdash \neg n$.
Since $\mathcal{S} \vdash n$,
the proposition is framework-variant in the sense relevant
to the neutrality requirements.
By Lemma~\ref{lem:contestability},
the substrate violates interpretive non-commitment
and extension stability.

The point is not that normative judgment is illegitimate
or irrelevant to accountability.
Rather, normative judgment belongs to interpretive frameworks.
If the substrate asserts that an action was permitted, prohibited,
obligatory, compliant, noncompliant, justified, or violative,
it embeds a framework-dependent conclusion
as a substrate-layer fact.
Any admissible framework that rejects that conclusion
cannot be layered atop the substrate without contradiction
or substrate revision.
Such a substrate no longer functions as a neutral substrate
for the relevant framework space.

%-------------------------------------------------------
\subsection{Causal Commitments}
\label{subsec:causal}
%-------------------------------------------------------

Causal attributions are likewise framework-dependent
in the sense relevant to Lemma~\ref{lem:contestability}.
Claims such as $\mathsf{Caused}(e_1,e_2)$
depend on background assumptions about causal mechanisms,
variable selection, counterfactual reasoning,
temporal scope, intervention conditions, and model boundaries.
Distinct causal frameworks may be admissible
while disagreeing about specific causal relationships.

By Assumption~\ref{asm:contestability},
if $\mathcal{S}$ commits to a causal proposition $c$,
then $c$ is not guaranteed invariant across admissible frameworks:
there exists some $\mathcal{F} \in \mathbb{F}$ such that
$\mathcal{S} \cup \mathcal{F} \vdash \neg c$.
Since $\mathcal{S} \vdash c$,
the proposition is framework-variant in the sense relevant
to the neutrality requirements.
By Lemma~\ref{lem:contestability},
the substrate violates
interpretive non-commitment and extension stability.

This argument does not rely on the presence
of an active dispute over a specific event.
It concerns the status of causal attribution
as a substrate-layer assertion.
As soon as a substrate asserts
$\mathsf{Caused}(e_1,e_2)$
as a foundational fact,
it commits to a particular causal interpretation
of the relation between $e_1$ and $e_2$.
An admissible framework that treats the two occurrences
as merely correlated,
as effects of a common cause,
or as causally related only under different model assumptions
cannot be layered atop the substrate without conflict.

To be clear, \textit{pre-causal does not mean acausal}.
The substrate does not deny that causation exists
nor does it deny that causal reasoning is necessary.
Rather, it refrains from embedding causal conclusions
as foundational commitments.
Causation is externalized to interpretive layers
where causal propositions are explicit,
contestable, and attributable.
Causation is externalized, not eliminated.

%-------------------------------------------------------
\subsection{Example: Incident Investigation}
\label{subsec:incident}
%-------------------------------------------------------

The following example illustrates the distinction between
substrate-layer referential structure,
framework-level conclusions,
and reified assertions about claims.

\begin{example}[Interpretive Disagreement in Incident Investigation]
  \label{ex:incident}

  Consider a substrate $\mathcal{S}$ representing a workplace incident.
  The substrate asserts only framework-invariant referential structure,
  including a referential regime for the relevant agents and occurrences.

  For example, the substrate may assert:

  \begin{itemize}
    \item $\mathsf{Agent}(a)$: $a$ is an agent;
    \item $\mathsf{Occurrence}(e_1)$ and $\mathsf{Occurrence}(e_2)$:
          $e_1$ and $e_2$ are occurrences;
    \item $\mathsf{Participant}(a,e_1)$:
          $a$ participated in $e_1$;
    \item $\mathsf{Before}(e_1,e_2)$:
          occurrence $e_1$ preceded occurrence $e_2$;
    \item referential-regime assertions fixing the individuation,
          co-reference, and persistence conditions for $a$, $e_1$, and $e_2$.
  \end{itemize}

  The substrate does not assert
  $\mathsf{Caused}(e_1,e_2)$,
  $\mathsf{Responsible}(a,e_2)$,
  or any normative conclusion about whether $a$ violated,
  satisfied, or breached an obligation.

  Two admissible interpretive frameworks extend $\mathcal{S}$:

  \medskip

  \noindent
  $\mathcal{F}_1$,
  an engineering analysis,
  applies a mechanistic causal model and derives
  $\mathsf{Caused}(e_1,e_2)$
  and
  $\mathsf{Responsible}(a,e_2)$.

  \noindent
  $\mathcal{F}_2$,
  a legal framework,
  applies jurisdiction-specific standards of responsibility
  and derives
  $\neg\mathsf{Responsible}(a,e_2)$.

  \medskip
  The two frameworks are mutually inconsistent with respect to
  the responsibility conclusion:
  $\mathcal{F}_1 \cup \mathcal{F}_2 \vdash \bot$.
  However, each is separately consistent with the substrate:
  $\mathcal{S} \cup \mathcal{F}_1 \not\vdash \bot$
  and
  $\mathcal{S} \cup \mathcal{F}_2 \not\vdash \bot$.
  The substrate remains neutral because it provides a stable referential base
  without asserting either framework's conclusion as a substrate-layer fact.

  \medskip
  \noindent
  \textbf{Violation case.}
  Suppose instead that $\mathcal{S}$ asserts
  $\mathsf{Responsible}(a,e_2)$
  as a substrate-layer fact.
  Then $\mathcal{S} \cup \mathcal{F}_2$
  entails both
  $\mathsf{Responsible}(a,e_2)$
  and
  $\neg\mathsf{Responsible}(a,e_2)$.
  Thus
  $\mathcal{S} \cup \mathcal{F}_2 \vdash \bot$,
  and the substrate violates extension stability.
  The failure is not that responsibility is irrelevant.
  The failure is that the substrate asserted a framework-level conclusion
  as a substrate-layer fact.

  Similarly, if $\mathcal{S}$ asserts
  $\mathsf{Caused}(e_1,e_2)$
  as a substrate-layer fact,
  then any admissible framework that rejects that causal attribution
  cannot extend the substrate without contradiction.
  The substrate has embedded a causal interpretation
  rather than merely providing the referential structure
  over which causal interpretations may be made.

  \medskip
  \noindent
  \textbf{Reification case.}
  Now suppose $\mathcal{S}$ neither asserts nor denies
  $\mathsf{Caused}(e_1,e_2)$
  or
  $\mathsf{Responsible}(a,e_2)$,
  but instead represents the engineering framework's conclusions
  as reified assertions:

  \[
    \mathsf{Asserts}(\mathcal{F}_1, \mathsf{Caused}(e_1,e_2))
  \]

  and

  \[
    \mathsf{Asserts}(\mathcal{F}_1, \mathsf{Responsible}(a,e_2))
  \]

  These are substrate-layer facts about the existence,
  source, and attribution of claims.
  They are not substrate-layer commitments to the truth
  of the claims' contents.
  Both $\mathcal{F}_1$ and $\mathcal{F}_2$
  can be layered atop $\mathcal{S}$ without contradiction:
  each framework reasons over the same referents
  while deriving its own conclusions.

  This is the structurally neutral pattern.
  The substrate records discourse about the world
  without asserting contested conclusions about the world.
\end{example}

%-------------------------------------------------------
\subsection{Reification and Contextual Modeling}
\label{subsec:reification}
%-------------------------------------------------------

The reification case in Example~\ref{ex:incident} illustrates
a broader point about contextual modeling approaches.
Representing causal or normative claims as reified assertions
attributed to agents,
institutions, documents, or frameworks,
is consistent with neutrality,
provided the substrate does not also assert the reified content
as a substrate-layer fact.

This pattern is compatible with DOLCE/DUL-style modeling of
situations and descriptions~\citep{borgo2022dolce}.
In such approaches, a description or situation may be represented
as an entity in the ontology,
with attribution, provenance, temporal scope,
and other metadata.
The substrate may assert that a description exists,
that it was made by a particular source,
or that it concerns particular referents.
It need not assert that the description is true.

Reification preserves interpretive non-commitment
only if the substrate refrains from simultaneously
asserting the reified content as a substrate-layer truth.
If a causal or normative relation is asserted
as a substrate-layer fact and also represented as a reified claim,
neutrality is already lost.
The reified representation does not cure
the substrate-layer assertion.
It merely duplicates it.

Thus, reification and contextual modeling do not circumvent
the theorem's constraint.
They satisfy it when they externalize interpretive commitments
rather than embedding them under a different representational form.
The constraint concerns what is asserted as foundational fact,
not the representational mechanism through which
interpretive discourse is recorded.

%-------------------------------------------------------
\subsection{The Ontological Neutrality Theorem}
\label{subsec:formal-proof}
%-------------------------------------------------------

The preceding observations yield the central result of this work.

\begin{theorem}[Ontological Neutrality Theorem]
  \label{thm:neutrality}
  Let $\mathcal{S}$ be a substrate ontology intended to function
  as a neutral substrate across admissible interpretive frameworks,
  and let $\mathbb{F}$ be the set of all admissible frameworks.
  Suppose the domain satisfies the contestability assumption
  for causal and normative propositions
  and the referential-invariance scope condition
  that the substrate's framework-invariant referential structure,
  including its referential regime,
  is not contested
  across the admissible frameworks under consideration.
  Then $\mathcal{S}$ satisfies the neutrality requirements
  of interpretive non-commitment and extension stability
  if and only if its foundational layer
  is restricted to framework-invariant referential structure and
  does not assert causal or normative conclusions as substrate-layer facts.
\end{theorem}

Here a \emph{causal or normative commitment} means the assertion
of a causal relation or normative proposition as a substrate-layer fact,
for example that one event caused another,
that an actor was responsible for an occurrence,
or that an action was permitted, prohibited,
obligatory, compliant, noncompliant, justified, or violative.
It does not include the representation of such propositions
as attributed, provenance-bearing, reified assertions.

\begin{proof}

  \noindent
  \textbf{Necessity.}
  Suppose $\mathcal{S}$ satisfies the neutrality requirements.
  Assume, for contradiction, that $\mathcal{S}$
  commits to a causal or normative proposition $p$.
  By Definition~\ref{def:substrate-layer-commitment},
  $\mathcal{S} \vdash p$.
  By Assumption~\ref{asm:contestability},
  there exists an admissible framework
  $\mathcal{F} \in \mathbb{F}$
  such that
  $\mathcal{S} \cup \mathcal{F} \vdash \neg p$.
  Since $\mathcal{S} \vdash p$,
  it follows that
  $\mathcal{S} \cup \mathcal{F} \vdash p$.
  Therefore
  $\mathcal{S} \cup \mathcal{F} \vdash p \land \neg p$,
  so
  $\mathcal{S} \cup \mathcal{F} \vdash \bot$.
  This contradicts extension stability.
  Moreover, because $p$ is rejected by an admissible framework
  while asserted by the substrate,
  $p$ is framework-variant in the relevant sense,
  so asserting it also violates interpretive non-commitment.
  Therefore, no neutral substrate may make
  causal or normative commitments.

  It remains to identify what may remain in the foundational layer.
  By the referential-invariance scope condition,
  the substrate may assert framework-invariant referential structure:
  the referential regimes by which entities, occurrences,
  and institutional artifacts are individuated,
  co-referred, and tracked across time,
  together with structural relations that assert neither causal efficacy
  nor normative force.
  Thus, under the stated scope condition,
  a neutral substrate's foundational layer must be restricted
  to framework-invariant referential structure
  and must exclude
  causal and normative commitments.

  \medskip
  \noindent
  \textbf{Sufficiency.}
  Suppose $\mathcal{S}$ is restricted at the foundational layer
  to framework-invariant referential structure
  and makes no causal or normative commitments.
  Let $\mathcal{F} \in \mathbb{F}$ be any admissible framework.

  Because the assertions of $\mathcal{S}$ are, by hypothesis,
  framework-invariant referential assertions, and because the
  framework-invariant referential structure $\mathcal{S}$ asserts,
  including its referential regime, is, by the scope condition,
  not contested across admissible frameworks, no assertion of
  $\mathcal{S}$ is contradicted by $\mathcal{F}$.
  Therefore
  $\mathcal{S} \cup \mathcal{F} \not\vdash \bot$.
  Since $\mathcal{F}$ was arbitrary,
  $\mathcal{S}$ satisfies extension stability.

  It remains to show interpretive non-commitment.
  Since $\mathcal{S}$ asserts only framework-invariant referential structure,
  there is no proposition $p$ such that
  $\mathcal{S} \vdash p$
  and $p$ is framework-variant with respect to $\mathbb{F}$.
  Causal and normative propositions may appear only as
  reified, attributed, provenance-bearing content.
  The substrate may assert that a framework made a claim;
  it does not assert the claim's truth.
  Thus $\mathcal{S}$ satisfies interpretive non-commitment.

  Therefore $\mathcal{S}$ satisfies both neutrality requirements:
  interpretive non-commitment and extension stability.

  \medskip
  \noindent
  \textit{Scope note.}
  The sufficiency direction depends on the referential-invariance
  scope condition.
  If admissible frameworks disagree about individuation,
  co-reference, or persistence---for example,
  whether two records denote the same entity,
  whether an occurrence is one event or many,
  or whether an institutional artifact persists across legal change---then
  the substrate cannot be neutral with respect to those frameworks
  until the referential contestation is resolved or scoped out.
  The theorem applies to domains in which the substrate's
  framework-invariant referential structure,
  including its referential regime,
  can be fixed prior to interpretive extension.
  It does not claim that all domains admit neutral substrates.
\end{proof}

This result is not contingent on a particular ontology,
syntax, protocol, or implementation technology.
It follows directly from the functional role
the substrate is required to play under persistent interpretive disagreement.
Neutrality across disagreement is structurally incompatible
with substrate-layer causal or normative commitments.

The next section considers the implications of this result for ontology design,
clarifying the permitted and excluded substrate-layer commitments
and the mechanisms by which interpretation, evaluation,
and explanation may be externalized without loss of accountability.
